\section{Задачи комбинаторной оптимизации}
\label{sec:combinatorial}

\noindent\hspace{1.25cm}\emph{Раздел используется в ЛР~№7 (часть Б).}

\subsection{Особенности применения генетических алгоритмов}

В задачах комбинаторной оптимизации решение~--- подмножество, перестановка или разбиение конечного
множества объектов, а число допустимых решений растёт экспоненциально или факториально с размером
задачи. Многие такие задачи NP-трудны, и ГА здесь~--- эвристика, находящая хорошие решения за
приемлемое время. Главные вопросы при его построении:
\begin{itemize}
\item \emph{представление}: кодирование должно покрывать все допустимые решения и по возможности только
их;
\item \emph{операторы}: скрещивание и мутация должны сохранять допустимость и передавать потомкам
существенные для качества свойства родителей (для маршрута~--- смежность городов, а не их абсолютные
позиции);
\item \emph{учёт ограничений}: недопустимые особи \emph{штрафуются} (штраф растёт с величиной нарушения),
\emph{восстанавливаются} (repair) эвристикой либо не возникают благодаря специальным операторам.
\end{itemize}

\subsection{Задачи о рюкзаке и о покрытии}

\emph{Задача о рюкзаке}: из $n$ предметов с весами $w_i$ и ценностями $c_i$ выбрать набор наибольшей
ценности, помещающийся в рюкзак вместимостью $W$:
\begin{equation}
\label{eq:knapsack}
\sum_{i=1}^{n} c_i x_i \to \max, \qquad \sum_{i=1}^{n} w_i x_i \le W, \qquad x_i \in \{0, 1\}.
\end{equation}
Двоичная хромосома $(x_1, \ldots, x_n)$ кодирует решение ($x_i = 1$~--- предмет взят). Перегрузку учитывают
штрафом $F = \sum c_i x_i - \lambda\max\bigl(0, \sum w_i x_i - W\bigr)$ или восстановлением: из рюкзака
удаляют предметы с наименьшим отношением $c_i/w_i$.

\emph{Задача о покрытии}: дана двоичная матрица $A = (a_{ij})$, $a_{ij} = 1$, если столбец $j$ стоимостью
$c_j$ покрывает строку $i$; требуется выбрать набор столбцов минимальной стоимости, покрывающий все строки:
\begin{equation}
\label{eq:cover}
\sum_{j=1}^{n} c_j x_j \to \min, \qquad \sum_{j=1}^{n} a_{ij} x_j \ge 1, \ i = 1, \ldots, m, \qquad
x_j \in \{0, 1\}.
\end{equation}
Хромосома~--- двоичный вектор выбранных столбцов; непокрытые строки учитываются штрафом или
восстановлением. К этой задаче сводятся размещение станций обслуживания и составление расписаний.

\subsection{Задача коммивояжёра}

Дано $n$ городов и матрица расстояний $d_{ij}$. Требуется найти замкнутый маршрут (тур),
проходящий через каждый город ровно один раз и имеющий минимальную длину:
\begin{equation}
\label{eq:tsp}
L(\pi) = \sum_{k=1}^{n-1} d_{\pi_k \pi_{k+1}} + d_{\pi_n \pi_1} \to \min_{\pi},
\end{equation}
где $\pi$~--- перестановка номеров городов. Для симметричной задачи число различных туров равно
$(n - 1)!/2$: для 76 городов это около $1,2 \cdot 10^{109}$, полный перебор невозможен. В тестовой
библиотеке TSPLIB~\cite{reinelt1991} для задач типа \code{EUC\_2D} расстояние~--- евклидово, округлённое
до ближайшего целого: $d_{ij} = \bigl\lfloor\sqrt{(x_i - x_j)^2 + (y_i - y_j)^2} + 0,5\bigr\rfloor$;
для многих задач библиотеки известны оптимальные туры (таблица~\ref{tab:l7-var}).

Обычные двоичные операторы порождают строки, не являющиеся перестановками, поэтому используются три
специальных представления тура~\cite{michalewicz1996, skobtsov2018}.

\emph{Представление соседства} (adjacency): в позиции $i$ записан город, следующий в туре за городом $i$.
Тур $1{-}2{-}4{-}3{-}8{-}5{-}9{-}6{-}7$ кодируется вектором $(2\ 4\ 8\ 3\ 9\ 7\ 1\ 5\ 6)$. Не всякий вектор
соответствует одному туру (возможны подциклы), поэтому применяются специальные операторы:
\emph{обмен рёбер} (потомок строится поочерёдным выбором ребра первого и второго родителей),
\emph{обмен подтуров} (чередуются фрагменты туров родителей) и \emph{эвристическое скрещивание}
(из текущего города выбирается более короткое из двух рёбер родителей). Представление явно работает с
рёбрами, которые и определяют длину тура.

\emph{Порядковое представление} (ordinal): задаётся упорядоченный список городов, например
$(1\ 2\ \ldots\ 9)$; очередной город тура кодируется номером его позиции в списке, после чего город
удаляется из списка. Тур $1{-}2{-}4{-}3{-}8{-}5{-}9{-}6{-}7$ кодируется как $(1\ 1\ 2\ 1\ 4\ 1\ 3\ 1\ 1)$.
Любой вектор с $i$-м геном из $[1;\ n - i + 1]$ является допустимым туром, поэтому работает обычный
одноточечный кроссинговер, но хвост хромосомы после разреза задаёт в потомке другие города, и хорошие
фрагменты не наследуются.

\emph{Представление путей} (path)~--- самое естественное: тур записывается как перестановка городов в
порядке обхода, $(1\ 2\ 4\ 3\ 8\ 5\ 9\ 6\ 7)$. Для него разработаны операторы, сохраняющие перестановку.
Рассмотрим их на родителях
\[
P_1 = (1\ 2\ 3\ |\ 4\ 5\ 6\ 7\ |\ 8\ 9), \qquad P_2 = (4\ 5\ 2\ |\ 1\ 8\ 7\ 6\ |\ 9\ 3)
\]
с точками разреза после третьей и седьмой позиций.

\emph{Частично соответствующий кроссинговер} (PMX, Д.~Голдберг и Р.~Лингл~\cite{goldberg1985}). Потомок
$C_1$ получает средний сегмент $P_1$: $(\cdot\ \cdot\ \cdot\ |\ 4\ 5\ 6\ 7\ |\ \cdot\ \cdot)$. Сегменты задают
отображение $4 \leftrightarrow 1$, $5 \leftrightarrow 8$, $6 \leftrightarrow 7$, $7 \leftrightarrow 6$.
Остальные позиции заполняются генами $P_2$; если ген уже есть в сегменте, он заменяется по отображению
(при необходимости многократно): $4 \to 1$, $5 \to 8$, $2$ и $9$, $3$ переносятся без изменений.
Получаем $C_1 = (1\ 8\ 2\ |\ 4\ 5\ 6\ 7\ |\ 9\ 3)$, аналогично $C_2 = (4\ 2\ 3\ |\ 1\ 8\ 7\ 6\ |\ 5\ 9)$. PMX
сохраняет абсолютные позиции большинства городов.

\emph{Упорядоченный кроссинговер} (OX, Л.~Дэвис). Потомок получает сегмент $P_1$, а свободные позиции,
начиная с позиции после второй точки разреза, заполняются городами $P_2$ в порядке их следования в
$P_2$ (также начиная после второй точки и по кругу), пропуская уже имеющиеся. Порядок обхода в $P_2$ с
девятой позиции: $9, 3, 4, 5, 2, 1, 8, 7, 6$; без городов сегмента $4, 5, 6, 7$ остаётся
$9, 3, 2, 1, 8$. Итог: $C_1 = (2\ 1\ 8\ |\ 4\ 5\ 6\ 7\ |\ 9\ 3)$, аналогично
$C_2 = (3\ 4\ 5\ |\ 1\ 8\ 7\ 6\ |\ 9\ 2)$. OX сохраняет относительный порядок городов, что для тура
важнее абсолютных позиций.

\emph{Циклический кроссинговер} (CX, И.~Оливер и др.~\cite{oliver1987}) не использует точек разреза. Позиции
разбиваются на \emph{циклы}: начиная с позиции~1, берём город $P_1[1] = 1$, смотрим, какой город стоит на
этой позиции в $P_2$ ($4$), находим позицию города $4$ в $P_1$ (позиция~4), и так далее, пока не
вернёмся к началу. Получаем циклы позиций $\{1, 4\}$, $\{2, 5, 8, 9, 3\}$ и $\{6, 7\}$
(рисунок~\ref{fig:cx}). Потомок $C_1$ берёт позиции первого цикла из $P_1$, второго~--- из $P_2$,
третьего~--- снова из $P_1$:
\[
C_1 = (1\ 5\ 2\ 4\ 8\ 6\ 7\ 9\ 3), \qquad C_2 = (4\ 2\ 3\ 1\ 5\ 7\ 6\ 8\ 9).
\]
Каждый город потомка стоит на той же позиции, что и у одного из родителей. Если родители состоят из
одного цикла, потомки совпадают с родителями.

\begin{figure}[!htb]
\centering
\begin{tikzpicture}
\def\A{cblue!22}\def\B{corange!30}\def\C{cgreen!35}
\bitrow[1.2]{0}{$P_1$}{1/\A,2/\B,3/\B,4/\A,5/\B,6/\C,7/\C,8/\B,9/\B}
\bitrow[1.2]{-0.6}{$P_2$}{4/\A,5/\B,2/\B,1/\A,8/\B,7/\C,6/\C,9/\B,3/\B}
\bitrow[1.2]{-1.6}{$C_1$}{1/\A,5/\B,2/\B,4/\A,8/\B,6/\C,7/\C,9/\B,3/\B}
\bitrow[1.2]{-2.2}{$C_2$}{4/\A,2/\B,3/\B,1/\A,5/\B,7/\C,6/\C,8/\B,9/\B}
\foreach \i in {1,...,9} \node[font=\scriptsize, text=cgray] at ($(1.2, 0.5) + (\i*\bitw - 0.3\bitw, 0)$) {\i};
\node[anchor=west, font=\small, align=left] at (6.0, -0.3) {цикл 1: позиции $\{1, 4\}$ \\
  цикл 2: позиции $\{2, 5, 8, 9, 3\}$ \\ цикл 3: позиции $\{6, 7\}$};
\node[anchor=west, font=\small, align=left] at (6.0, -1.9) {$C_1$: циклы 1, 3 из $P_1$, цикл 2 из $P_2$ \\
  $C_2$: циклы 1, 3 из $P_2$, цикл 2 из $P_1$};
\end{tikzpicture}
\caption{Циклический кроссинговер: цветом выделены циклы позиций}
\label{fig:cx}
\end{figure}

\emph{Скрещивание с рекомбинацией рёбер} (edge recombination, ER, Д.~Уитли~\cite{whitley1989}) строит
для каждого города список его соседей в обоих родителях и формирует потомка, переходя в соседний город
с наименьшим числом оставшихся соседей. Почти все рёбра потомка унаследованы от родителей, поэтому ER
обычно эффективнее PMX, OX и CX, сохраняющих позиции или порядок, но не смежность городов.

\emph{Мутации перестановок}: \emph{обмен} (swap)~--- два города меняются местами (изменяются четыре
ребра); \emph{вставка}~--- город переносится на другую позицию (три ребра); \emph{инверсия}~--- отрезок
тура записывается в обратном порядке (для симметричной задачи изменяются только два ребра);
\emph{перемешивание} случайного отрезка.

\subsection{Локальный поиск 2-opt и меметические алгоритмы}

Операторы, сохраняющие позиции (CX, PMX), плохо передают потомкам главное свойство хорошего тура~---
набор коротких рёбер. Поэтому чистый ГА для задачи коммивояжёра сходится медленно. Существенно
лучшие результаты даёт сочетание ГА с \emph{локальным поиском}.

\emph{Алгоритм 2-opt} (Г.~Крос, 1958~\cite{croes1958}) удаляет из тура два ребра $(a, b)$ и $(c, d)$ и
соединяет концы крест-накрест рёбрами $(a, c)$ и $(b, d)$, что эквивалентно инверсии отрезка тура от $b$
до $c$ (рисунок~\ref{fig:2opt}). Замена выполняется, если она укорачивает тур:
\begin{equation}
\label{eq:2opt}
\Delta = d_{ac} + d_{bd} - d_{ab} - d_{cd} < 0 .
\end{equation}
Перебор всех $O(n^2)$ пар рёбер повторяется, пока есть улучшающая замена. Результат~---
\emph{2-оптимальный} тур без самопересечений, в среднем примерно на 5\,\% длиннее оптимального (более
сильные эвристики~--- 3-opt и алгоритм Лина~--- Кернигана).

\begin{figure}[!htb]
\centering
\begin{tikzpicture}[city/.style={circle, fill=black, inner sep=1.6pt}, >=Stealth]
\foreach \dx/\after in {0/0, 7.8/1} {
\begin{scope}[yscale=0.8]
\begin{scope}[xshift=\dx cm]
\node[city, label=below left:$a$] (a) at (0, 0) {};
\node[city, label=above right:$b$] (b) at (3, 2) {};
\node[city, label=above left:$c$] (c) at (0, 2) {};
\node[city, label=below right:$d$] (d) at (3, 0) {};
\draw[thick, cblue, dashed] (b) .. controls (2.4, 3.1) and (0.6, 3.1) .. (c);
\draw[thick, cblue, dashed] (d) .. controls (2.4, -1.1) and (0.6, -1.1) .. (a);
\ifnum\after=0
  \draw[thick, corange] (a) -- (b);
  \draw[thick, corange] (c) -- (d);
\else
  \draw[thick, cgreen] (a) -- (c);
  \draw[thick, cgreen] (b) -- (d);
\fi
\end{scope}
\end{scope}}
\node[font=\small] at (1.5, -1.45) {до: рёбра $(a, b)$ и $(c, d)$ пересекаются};
\node[font=\small] at (9.3, -1.45) {после: рёбра $(a, c)$ и $(b, d)$};
\draw[->, thick, cgray] (4.6, 1) -- (6.4, 1);
\end{tikzpicture}
\caption{Ход 2-opt: рёбра $(a, b)$, $(c, d)$ заменяются на $(a, c)$, $(b, d)$, и путь от $b$ до $c$
проходится в обратном направлении (пунктир~--- остальная часть тура)}
\label{fig:2opt}
\end{figure}

\emph{Меметический алгоритм} (П.~Москато, 1989)~--- ГА, в котором каждый потомок улучшается локальным
поиском до ближайшего локального оптимума: ГА комбинирует фрагменты разных локальных оптимумов, а
локальный поиск быстро доводит решения. На наборе eil76 меметический ГА находит оптимальный тур, а
ГА с CX без локального поиска отстаёт от оптимума на 40--60\,\% (см. пример выполнения ЛР~№7).

\label{sec:py-tsp}\emph{Реализация на Python.} Тур хранится как перестановка индексов городов, популяция~--- массивом $(N, n)$; длина всех туров
популяции вычисляется одной строкой \code{D[pop, np.roll(pop, -1, axis=1)].sum(axis=1)}. Выигрыш 2-opt
(\ref{eq:2opt}) для всех пар рёбер считается одной операцией над массивами индексов, и случайный тур из
76 городов доводится до 2-оптимального за несколько миллисекунд (приложение~\ref{app:listings}).
PMX и OX для проверки есть в DEAP (\code{tools.cxPartialyMatched}, \code{tools.cxOrdered}).
